%! The Secant, Cosecant, and Cotangent Functions — Unit 3, Lesson 3.11 — Student Packet Cover Sheet
\documentclass[10pt]{article}
\usepackage{precalculus-article}
\usepackage{precalculus-boxes}
\usepackage{ltablex}
\keepXColumns

\begin{document}

% Full-bleed navy banner
\begin{tikzpicture}[remember picture, overlay]
  \fill[navy] (current page.north west)
    rectangle ([yshift=-0.9in]current page.north east);
\end{tikzpicture}
\vspace{-0.8in}

\begin{center}
  {\color{white}\textbf{\LARGE Precalculus}}\\[4pt]
  {\color{skymid}\large\textbf{Unit 3: Trigonometric and Polar Functions}}\\[6pt]
  {\color{white}\textbf{\large Lesson 3.11 \quad The Secant, Cosecant, and Cotangent Functions}}
\end{center}

\vspace{0.22in}
\namedateperiod
\vspace{0.18in}

\begin{learningtargetbox}
\small
\begin{enumerate}[leftmargin=1.5em, itemsep=5pt]
  \item I can define secant, cosecant, and cotangent as reciprocals of cosine,
    sine, and tangent, and evaluate each at unit-circle angles.
    \hfill\textit{(LO 3.11.A)}
  \item I can identify the domain restrictions and vertical asymptote locations
    for $\sec\theta$, $\csc\theta$, and $\cot\theta$.
    \hfill\textit{(LO 3.11.A, EK 3.11.A.3)}
  \item I can state the range of the secant and cosecant functions and explain
    why neither function produces output values between $-1$ and $1$.
    \hfill\textit{(LO 3.11.A, EK 3.11.A.3)}
\end{enumerate}
\end{learningtargetbox}

\vspace{0.14in}

\begin{tocbox}
\small
\renewcommand{\arraystretch}{1.5}
\begin{tabularx}{\linewidth}{@{} c l X r @{}}
\rowcolor{sky}
\textbf{\#} & \textbf{Component} & \textbf{Description} & \textbf{Score} \\
\midrule
1 & Warm-Up        & Spiral review: unit-circle cos/sin values, tangent as a quotient, asymptote behavior, reciprocal arithmetic & \blank{1.2cm} \\
2 & Guided Notes   & Definitions of sec, csc, cot; value tables; domain, range, asymptotes; summary of all three & \blank{1.2cm} \\
3 & Group Activity & Compute reciprocal trig values; explore domain restrictions and asymptotes (tiered R/A/E) & \blank{1.2cm} \\
4 & Exit Ticket    & Given a cosine value, find secant; state domain of csc; identify period of cot & \blank{1.2cm} \\
5 & Homework       & Evaluate sec/csc/cot; identify domains and asymptotes; connect sec graph to cos zeros; AP-style MC & \blank{1.2cm} \\
\midrule
  & \multicolumn{2}{r}{\textbf{Total}} & \blank{1.2cm} \\
\end{tabularx}
\end{tocbox}

\end{document}
